%%%%%%%%%%%%%%%%%%%%%%%%%%%%%
% Dr. Sette Diop            %
% Charge de Recherche CNRS  %
% LAGEP,  UCB Lyon I        %
% 43 Bd. 11 Nov. 1918       %
% 69622 Villeurbanne cedex  %
% France                    %
%                           %
% Tel. +33 72448516         %
% Fax. +33 72431197         %
%%%%%%%%%%%%%%%%%%%%%%%%%%%%%
%%%%%%%%%%%%%%%%%%%%%%%%%%%%%%%%%%%%%%%%%%%%%%%%%%%%%%%%%%%%%%%%%%%%%%%%%%%%
%%%%%%%%%%%%%%%%%%%%%% This is an input file for TeX. %%%%%%%%%%%%%%%%%%%%%%
%%%%%%%%%%%%%%%%%%%%%%%%%%%%%%%%%%%%%%%%%%%%%%%%%%%%%%%%%%%%%%%%%%%%%%%%%%%%

%%%%%%%%%%%%%%%%%%%%%%%%%%%%%%%%%%%%%%%%%%%%%%%%%%%%%%%%%%%%%%%%%%%%%%%%%%%%
%                                   Fonts                                  %
%%%%%%%%%%%%%%%%%%%%%%%%%%%%%%%%%%%%%%%%%%%%%%%%%%%%%%%%%%%%%%%%%%%%%%%%%%%%

\font\teneufb=eufb10
\font\seveneufb=eufb7
\font\fiveeufb=eufb5
\newfam\eufbfam
\textfont\eufbfam=\teneufb
\scriptfont\eufbfam=\seveneufb
\scriptscriptfont\eufbfam=\fiveeufb
\def\eufb#1{{\fam\eufbfam\relax#1}}

\font\teneufm=eufm10
\font\seveneufm=eufm7
\font\fiveeufm=eufm5
\newfam\eufmfam
\textfont\eufmfam=\teneufm
\scriptfont\eufmfam=\seveneufm
\scriptscriptfont\eufmfam=\fiveeufm
\def\eufm#1{{\fam\eufmfam\relax#1}}

\font\teneurm=eurm10
\font\seveneurm=eurm7
\font\fiveeurm=eurm5
\newfam\eurmfam
\textfont\eurmfam=\teneurm
\scriptfont\eurmfam=\seveneurm
\scriptscriptfont\eurmfam=\fiveeurm
\def\eurm#1{{\fam\eurmfam\relax#1}}

\font\teneusb=eusb10
\font\seveneusb=eusb7
\font\fiveeusb=eusb5
\newfam\eusbfam
\textfont\eusbfam=\teneusb
\scriptfont\eusbfam=\seveneusb
\scriptscriptfont\eusbfam=\fiveeusb
\def\eusb#1{{\fam\eusbfam\relax#1}}

\font\teneusm=eusm10
\font\seveneusm=eusm7
\font\fiveeusm=eusm5
\newfam\eusmfam
\textfont\eusmfam=\teneusm
\scriptfont\eusmfam=\seveneusm
\scriptscriptfont\eusmfam=\fiveeusm
\def\eusm#1{{\fam\eusmfam\relax#1}}

\font \tenbsy=cmbsy10
% \font \ninebsy=cmbsy9
% \font \eightbsy=cmbsy8
% \font \sevenbsy=cmbsy7
% \font \sixbsy=cmbsy6
% \font \fivebsy=cmbsy5

% Computer Modern Bold Extended
\font \twentyfourbf = cmbx10 scaled \magstep5
\font \twentybf = cmbx10 scaled \magstep4
\font \seventeenbf = cmbx10 scaled \magstep3
\font \fourteenbf = cmbx10 scaled \magstep2
\font \twelvebf = cmbx10 scaled \magstep1
\font \elevenbf = cmbx10 scaled \magstephalf
\font \tenbf = cmbx10
\font \ninebf = cmbx9
\font \eightbf = cmbx8
\font \sevenbf = cmbx7
\font \sixbf = cmbx6
\font \fivebf = cmbx5

\font \twentyfourbxti = cmbxti10 scaled \magstep5
\font \twentybxti = cmbxti10 scaled \magstep4
\font \seventeenbxti = cmbxti10 scaled \magstep3
\font \fourteenbxti = cmbxti10 scaled \magstep2
\font \twelvebxti = cmbxti10 scaled \magstep1
\font \elevenbxti = cmbxti10 scaled \magstephalf
\font \tenbxti = cmbxti10
% \font \ninebxti=cmbxti9    % Lacking
% \font \eightbxti=cmbxti8   % Lacking
% \font \sevenbxti=cmbxti7   % Lacking
% \font \sixbxti=cmbxti6     % Lacking
% \font \fivebxti=cmbxti5    % Lacking

% \font \twentyfoursc=cmcsc10 scaled \magstep5
% \font \twentysc=cmcsc10 scaled \magstep4
% \font \seventeensc=cmcsc10 scaled \magstep3
% \font \fourteensc=cmcsc10 scaled \magstep2
% \font \twelvesc=cmcsc10 scaled \magstep1
% \font \elevencsc=cmcsc10 scaled \magstephalf
% \font \tencsc=cmcsc10
% \font \ninecsc=cmcsc9
% \font \eightcsc=cmcsc8
% \font \sevencsc=cmcsc7   % Lacking
% \font \sixcsc=cmcsc6     % Lacking
% \font \fivecsc=cmcsc5    % Lacking

% Computer Modern Math Extension
\font \twentyfourex=cmex10 scaled \magstep5
\font \twentyex=cmex10 scaled \magstep4
\font \seventeenex=cmex10 scaled \magstep3
\font \fourteenex=cmex10 scaled \magstep2
\font \twelveex=cmex10 scaled \magstep1
\font \elevenex=cmex10 scaled \magstephalf
\font \tenex=cmex10
% \font \nineex=cmex9
% \font \eightex=cmex8
% \font \sevenex=cmex7
% \font \sixex=cmex6    % Lacking
% \font \fiveex=cmex5   % Lacking

% CM Funny Roman
\font \twentyfourfi = cmfi10 scaled \magstep5
\font \twentyfi = cmfi10 scaled \magstep4
\font \seventeenfi = cmfi10 scaled \magstep3
\font \fourteenfi = cmfi10 scaled \magstep2
\font \twelvefi = cmfi10 scaled \magstep1
\font \elevenfi = cmfi10 scaled \magstephalf
\font \tenfi = cmfi10
% \font \ninefi = cmfi9    % Lacking
% \font \eightfi = cmfi8   % Lacking
% \font \sevenfi = cmfi7   % Lacking
% \font \sixfi = cmfi6     % Lacking
% \font \fivefi = cmfi5    % Lacking

% \font \tenitt=cmitt10
% \font \nineitt=cmitt9    % Lacking
% \font \eightitt=cmitt8   % Lacking
% \font \sevenitt=cmitt7   % Lacking
% \font \sixitt=cmitt6     % Lacking
% \font \fiveitt=cmitt5    % Lacking

% Computer Modern Math Italic
\font \twentyfouri=cmmi10 scaled \magstep5
\font \twentyi=cmmi10 scaled \magstep4
\font \seventeeni=cmmi10 scaled \magstep3
\font \fourteeni=cmmi10 scaled \magstep2
\font \twelvei=cmmi10 scaled \magstep1
\font \eleveni=cmmi10 scaled \magstephalf
\font \teni=cmmi10
\font \ninei=cmmi9
\font \eighti=cmmi8
\font \seveni=cmmi7
\font \sixi=cmmi6
\font \fivei=cmmi5

\font \tenib=cmmib10
% \font \nineib=cmmib9
% \font \eightib=cmmib8
\font \sevenib=cmmib7
% \font \sixib=cmmib6
\font \fiveib=cmmib5

% Computer Modern Text Roman
\font \twentyfourrm=cmr10 scaled \magstep5
\font \twentyrm=cmr10 scaled \magstep4
\font \seventeenrm=cmr10 scaled \magstep3
\font \fourteenrm=cmr10 scaled \magstep2
\font \twelverm=cmr10 scaled \magstep1
\font \elevenrm=cmr10 scaled \magstephalf
\font \tenrm=cmr10
\font \ninerm=cmr9
\font \eightrm=cmr8
\font \sevenrm=cmr7
\font \sixrm=cmr6
\font \fiverm=cmr5

% Computer Modern Slanted Roman
\font \twentyfoursl=cmsl10 scaled \magstep5
\font \twentysl=cmsl10 scaled \magstep4
\font \seventeensl=cmsl10 scaled \magstep3
\font \fourteensl=cmsl10 scaled \magstep2
\font \twelvesl=cmsl10 scaled \magstep1
\font \elevensl=cmsl10 scaled \magstephalf
\font \tensl=cmsl10
\font \ninesl=cmsl9
\font \eightsl=cmsl8
% \font \sevensl=cmsl7   % Lacking
% \font \sixsl=cmsl6     % Lacking
% \font \fivesl=cmsl5    % Lacking

% \font \tensltt=cmsltt10
% \font \ninesltt=cmsltt9    % Lacking
% \font \eightsltt=cmsltt8   % Lacking
% \font \sevensltt=cmsltt7   % Lacking
% \font \sixsltt=cmsltt6     % Lacking
% \font \fivesltt=cmsltt5    % Lacking

\font \tenss=cmss10
\font \niness=cmss9
\font \eightss=cmss8
% \font \sevenss=cmss7      % Lacking
% \font \sixss=cmss6        % Lacking
% \font \fivess=cmss5       % Lacking

% \font \tenssb=cmssbx10
% \font \ninessb=cmssbx9    % Lacking
% \font \eightssb=cmssbx8   % Lacking
% \font \sevenssb=cmssbx7   % Lacking
% \font \sixssb=cmssbx6     % Lacking
% \font \fivessb=cmssbx5    % Lacking

% \font \tenssd=cmssdc10
% \font \ninessd=cmssdc9    % Lacking
% \font \eightssd=cmssdc8   % Lacking
% \font \sevenssd=cmssdc7   % Lacking
% \font \sixssd=cmssdc6     % Lacking
% \font \fivessd=cmssdc5    % Lacking

\font \tenssi=cmssi10
\font \ninessi=cmssi9
\font \eightssi=cmssi8
% \font \sevenssi=cmssi7   % Lacking
% \font \sixssi=cmssi6     % Lacking
% \font \fivessi=cmssi5    % Lacking

% Computer Modern Math Symbols
\font \twentyfoursy=cmsy10 scaled \magstep5
\font \twentysy=cmsy10 scaled \magstep4
\font \seventeensy=cmsy10 scaled \magstep3
\font \fourteensy=cmsy10 scaled \magstep2
\font \twelvesy=cmsy10 scaled \magstep1
\font \elevensy=cmsy10 scaled \magstephalf
\font \tensy=cmsy10
\font \ninesy=cmsy9
\font \eightsy=cmsy8
\font \sevensy=cmsy7
\font \sixsy=cmsy6
\font \fivesy=cmsy5

% Computer Modern Text Italic
\font \twentyfourit=cmti10 scaled \magstep5
\font \twentyit=cmti10 scaled \magstep4
\font \seventeenit=cmti10 scaled \magstep3
\font \fourteenit=cmti10 scaled \magstep2
\font \twelveit=cmti10 scaled \magstep1
\font \elevenit=cmti10 scaled \magstephalf
\font \tenit=cmti10
\font \nineit=cmti9
\font \eightit=cmti8
\font \sevenit=cmti7
% \font \sixit=cmti6    % Lacking
% \font \fiveit=cmti5   % Lacking

% Computer Modern True Type
\font \twentyfourtt=cmtt10 scaled \magstep5
\font \twentytt=cmtt10 scaled \magstep4
\font \seventeentt=cmtt10 scaled \magstep3
\font \fourteentt=cmtt10 scaled \magstep2
\font \twelvett=cmtt10 scaled \magstep1
\font \eleventt=cmtt10 scaled \magstephalf
\font \tentt=cmtt10
\font \ninett=cmtt9
\font \eighttt=cmtt8
% \font \seventt=cmtt7   % Lacking
% \font \sixtt=cmtt6     % Lacking
% \font \fivett=cmtt5    % Lacking

% \font \tenvt=cmvtt10
% \font \ninevt=cmvtt9    % Lacking
% \font \eightvt=cmvtt8   % Lacking
% \font \sevenvt=cmvtt7   % Lacking
% \font \sixvt=cmvtt6     % Lacking
% \font \fivevt=cmvtt5    % Lacking

% AMS Fonts Math Symbols (Group 2)
\font \twentyfourbm = msbm10 scaled \magstep5
\font \twentybm = msbm10 scaled \magstep4
\font \seventeenbm = msbm10 scaled \magstep3
\font \fourteenbm = msbm10 scaled \magstep2
\font \twelvebm = msbm10 scaled \magstep1
\font \elevenbm = msbm10 scaled \magstephalf
\font \tenbm = msbm10
% \font \ninebm = msbm9      % Lacking
% \font \eightbm = msbm8     % Lacking
\font \sevenbm = msbm7
% \font \sixbm = msbm6     % Lacking
\font \fivebm = msbm5

%%%%%%%%%%%%%%%%%%%%%%%%%%%%%%%%%%%%%%%%%%%%%%%%%%%%%%%%%%%%%%%%%%%%%%%%%%%%
%                    Size-switching font facilities                        %
%%%%%%%%%%%%%%%%%%%%%%%%%%%%%%%%%%%%%%%%%%%%%%%%%%%%%%%%%%%%%%%%%%%%%%%%%%%%

\newskip\ttglue
\def\twentypoint{\def\rm{\fam0\twentyrm} % switch to 20-point type
  \textfont0 = \twentyrm \scriptfont0 = \fourteenrm \scriptscriptfont0 = \tenrm
  \textfont1 = \twentyi  \scriptfont1 = \fourteeni  \scriptscriptfont1 = \teni
  \textfont2 = \twentysy \scriptfont2 = \fourteensy \scriptscriptfont2 = \tensy
  \textfont3 = \twentyex \scriptfont3 = \twentyex   \scriptscriptfont3 =
\twentyex
  \def\it{\fam\itfam\twentyit}
  \textfont\itfam = \twentyit
  \def\sl{\fam\slfam\twentysl}
  \textfont\slfam = \twentysl
  \def\tt{\fam\ttfam\twentytt}
  \textfont\ttfam = \twentytt
  \def\bf{\fam\bffam\twentybf}
  \textfont\bffam = \twentybf
  \scriptfont\bffam = \fourteenbf
  \scriptscriptfont\bffam = \tenbf
%  \def\sc{\fam\scfam\twentycsc}
%  \textfont\scfam = \twentycsc
  \tt \ttglue = .5em plus.25em minus.15em
  \normalbaselineskip = 25pt
  \setbox\strutbox = \hbox{\vrule height20pt depth5pt width0pt}
  \normalbaselines\rm}

\def\seventeenpoint{\def\rm{\fam0\seventeenrm} % switch to 17-point type
  \textfont0 = \seventeenrm \scriptfont0 = \twelverm    \scriptscriptfont0
= \ninerm
  \textfont1 = \seventeeni  \scriptfont1 = \twelvei     \scriptscriptfont1
= \ninei
  \textfont2 = \seventeensy \scriptfont2 = \twelvesy    \scriptscriptfont2
= \ninesy
  \textfont3 = \seventeenex \scriptfont3 = \seventeenex \scriptscriptfont3
= \seventeenex
  \def\it{\fam\itfam\seventeenit}
  \textfont\itfam = \seventeenit
  \def\sl{\fam\slfam\seventeensl}
  \textfont\slfam = \seventeensl
  \def\tt{\fam\ttfam\seventeentt}
  \textfont\ttfam = \seventeentt
  \def\bf{\fam\bffam\seventeenbf}
  \textfont\bffam = \seventeenbf
  \scriptfont\bffam = \twelvebf
  \scriptscriptfont\bffam = \ninebf
%  \def\sc{\fam\scfam\seventeencsc}
%  \textfont\scfam = \seventeencsc
  \tt \ttglue = .5em plus.25em minus.15em
  \normalbaselineskip = 21pt
  \setbox\strutbox = \hbox{\vrule height17pt depth4pt width0pt}
  \normalbaselines\rm}

\def\fourteenpoint{\def\rm{\fam0\fourteenrm} % switch to 14-point type
  \textfont0 = \fourteenrm \scriptfont0 = \tenrm      \scriptscriptfont0 =
\sevenrm
  \textfont1 = \fourteeni  \scriptfont1 = \teni       \scriptscriptfont1 =
\seveni
  \textfont2 = \fourteensy \scriptfont2 = \tensy      \scriptscriptfont2 =
\sevensy
  \textfont3 = \fourteenex \scriptfont3 = \fourteenex \scriptscriptfont3 =
\fourteenex
  \def\it{\fam\itfam\fourteenit}
  \textfont\itfam = \fourteenit
  \def\sl{\fam\slfam\fourteensl}
  \textfont\slfam = \fourteensl
  \def\tt{\fam\ttfam\fourteentt}
  \textfont\ttfam = \fourteentt
  \def\bf{\fam\bffam\fourteenbf}
  \textfont\bffam = \fourteenbf
  \scriptfont\bffam = \tenbf
  \scriptscriptfont\bffam = \sevenbf
%  \def\sc{\fam\scfam\fourteencsc}
%  \textfont\scfam = \fourteencsc
  \tt \ttglue = .5em plus.25em minus.15em
  \normalbaselineskip = 17pt
  \setbox\strutbox = \hbox{\vrule height11.9pt depth6.3pt width0pt}
  \normalbaselines\rm}

\def\twelvepoint{\def\rm{\fam0\twelverm} % switch to 12-point type
  \textfont0 = \twelverm \scriptfont0 = \ninerm   \scriptscriptfont0 = \sevenrm
  \textfont1 = \twelvei  \scriptfont1 = \ninei    \scriptscriptfont1 = \seveni
  \textfont2 = \twelvesy \scriptfont2 = \ninesy   \scriptscriptfont2 = \sevensy
  \textfont3 = \twelveex \scriptfont3 = \twelveex \scriptscriptfont3 = \twelveex
  \def\it{\fam\itfam\twelveit}
  \textfont\itfam = \twelveit
  \def\sl{\fam\slfam\twelvesl}
  \textfont\slfam = \twelvesl
  \def\tt{\fam\ttfam\twelvett}
  \textfont\ttfam = \twelvett
  \def\bf{\fam\bffam\twelvebf}
  \textfont\bffam = \twelvebf
  \scriptfont\bffam = \ninebf
  \scriptscriptfont\bffam = \sevenbf
%  \def\sc{\fam\scfam\twelvecsc}
%  \textfont\scfam = \twelvecsc
  \tt \ttglue = .5em plus.25em minus.15em
  \normalbaselineskip = 14pt
  \setbox\strutbox = \hbox{\vrule height9.5pt depth4.5pt width0pt}
  \normalbaselines\rm}

\def\elevenpoint{\def\rm{\fam0\elevenrm} % switch to 11-point type
  \textfont0 = \elevenrm \scriptfont0 = \eightrm  \scriptscriptfont0 = \sixrm
  \textfont1 = \eleveni  \scriptfont1 = \eighti   \scriptscriptfont1 = \sixi
  \textfont2 = \elevensy \scriptfont2 = \eightsy  \scriptscriptfont2 = \sixsy
  \textfont3 = \elevenex \scriptfont3 = \elevenex \scriptscriptfont3 = \elevenex
  \def\it{\fam\itfam\elevenit}
  \textfont\itfam = \elevenit
  \def\sl{\fam\slfam\elevensl}
  \textfont\slfam = \elevensl
  \def\tt{\fam\ttfam\eleventt}
  \textfont\ttfam = \eleventt
  \def\bf{\fam\bffam\elevenbf}
  \textfont\bffam = \elevenbf
  \scriptfont\bffam = \eightbf
  \scriptscriptfont\bffam = \sixbf
%  \def\sc{\fam\scfam\elevencsc}
%  \textfont\scfam = \elevencsc
  \tt \ttglue = .5em plus.25em minus.15em
  \normalbaselineskip = 14pt
  \setbox\strutbox = \hbox{\vrule height9pt depth4pt width0pt}
  \normalbaselines\rm}

\def\tenpoint{\def\rm{\fam0\tenrm} %  switch to 10-point type
  \textfont0 = \tenrm \scriptfont0 = \sevenrm \scriptscriptfont0 = \fiverm
  \textfont1 = \teni  \scriptfont1 = \seveni  \scriptscriptfont1 = \fivei
  \textfont2 = \tensy \scriptfont2 = \sevensy \scriptscriptfont2 = \fivesy
  \textfont3 = \tenex \scriptfont3 = \tenex   \scriptscriptfont3 = \tenex
  \def\it{\fam\itfam\tenit}
  \textfont\itfam = \tenit
  \def\sl{\fam\slfam\tensl}
  \textfont\slfam=\tensl
  \def\tt{\fam\ttfam\tentt}
  \textfont\ttfam = \tentt
  \def\bf{\fam\bffam\tenbf}
  \textfont\bffam = \tenbf
  \scriptfont\bffam = \sevenbf
  \scriptscriptfont\bffam = \fivebf
  \tt \ttglue = .5em plus.25em minus.15em
  \normalbaselineskip = 12pt
  \setbox\strutbox=\hbox{\vrule height8.5pt depth3.5pt width0pt}
  \let\sc = \eightrm \let\big = \tenbig \normalbaselines\rm}

\def\ninepoint{\def\rm{\fam0\ninerm} % switch to 9-point type
  \textfont0 = \ninerm \scriptfont0 = \sixrm \scriptscriptfont0 = \fiverm
  \textfont1 = \ninei  \scriptfont1 = \sixi  \scriptscriptfont1 = \fivei
  \textfont2 =\ninesy  \scriptfont2 = \sixsy \scriptscriptfont2 = \fivesy
  \textfont3 = \tenex  \scriptfont3 = \tenex \scriptscriptfont3 = \tenex
  \def\it{\fam\itfam\nineit}
  \textfont\itfam = \nineit
  \def\sl{\fam\slfam\ninesl}
  \textfont\slfam = \ninesl
  \def\tt{\fam\ttfam\ninett}
  \textfont\ttfam = \ninett
  \def\bf{\fam\bffam\ninebf}
  \textfont\bffam = \ninebf
  \scriptfont\bffam = \sixbf
  \scriptscriptfont\bffam = \fivebf
  \tt \ttglue = .5em plus.25em minus.15em
  \normalbaselineskip = 11pt
  \setbox\strutbox = \hbox{\vrule height8pt depth3pt width0pt}
  \let\sc = \sevenrm \let\big = \ninebig \normalbaselines\rm}

\def\eightpoint{\def\rm{\fam0\eightrm} % switch to 8-point type
  \textfont0 = \eightrm \scriptfont0 = \sixrm \scriptscriptfont0 = \fiverm
  \textfont1 = \eighti  \scriptfont1 = \sixi  \scriptscriptfont1 = \fivei
  \textfont2 = \eightsy \scriptfont2 = \sixsy \scriptscriptfont2 = \fivesy
  \textfont3 = \tenex \scriptfont3 = \tenex \scriptscriptfont3 = \tenex
  \def\it{\fam\itfam\eightit}
  \textfont\itfam = \eightit
  \def\sl{\fam\slfam\eightsl}
  \textfont\slfam = \eightsl
  \def\tt{\fam\ttfam\eighttt}
  \textfont\ttfam = \eighttt
  \def\bf{\fam\bffam\eightbf}
  \textfont\bffam = \eightbf
  \scriptfont\bffam = \sixbf
  \scriptscriptfont\bffam = \fivebf
  \tt \ttglue = .5em plus.25em minus.15em
  \normalbaselineskip = 9pt
  \setbox\strutbox=\hbox{\vrule height7pt depth2pt width0pt}
  \let\sc = \sixrm \let\big = \eightbig \normalbaselines\rm}
\def\sevenpoint{\def\rm{\fam0\sevenrm} % switch to a 7-point type
  \textfont0 = \sevenrm \scriptfont0 = \fiverm \scriptscriptfont0 = \fiverm
  \textfont1 = \seveni  \scriptfont1 = \fivei  \scriptscriptfont1 = \fivei 
  \textfont2 = \sevensy \scriptfont2 = \fivesy \scriptscriptfont2 = \fivesy
  \textfont3 = \tenex   \scriptfont3 = \tenex  \scriptscriptfont3 = \tenex 
  \def\it{\fam\itfam\sevenit}
  \textfont\itfam = \sevenit
  \def\sl{\fam\slfam\sevenit} % \sevensl is lacking. Replaced by \sevenit
  \textfont\slfam = \sevenit  % \sevensl is lacking. Replaced by \sevenit
  \def\tt{\fam\ttfam\sevenit} % \seventt is lacking. Replaced by \sevenit
  \textfont\ttfam = \sevenit  % \seventt is lacking. Replaced by \sevenit
  \def\bf{\fam\bffam\sevenbf}
  \textfont\bffam = \sevenbf
  \scriptfont\bffam = \fivebf
  \scriptscriptfont\bffam = \fivebf
  \tt \ttglue = .5em plus.25em minus.15em
  \normalbaselineskip = 10pt
  \setbox\strutbox = \hbox{\vrule height7.5pt depth2.5pt width0pt}
  \let\sc = \sixrm \let\big = \sevenbig \normalbaselines\rm}

\def\tenbig#1{{\hbox{$\left#1\vbox to8.5pt{}\right.\n@space$}}}
\def\ninebig#1{{\hbox{$\textfont0 = \tenrm\textfont2 = \tensy
  \left#1\vbox to7.25pt{}\right.\n@space$}}}
\def\eightbig#1{{\hbox{$\textfont0 = \ninerm\textfont2 = \ninesy
  \left#1\vbox to6.5pt{}\right.\n@space$}}}
\def\sevenbig#1{{\hbox{$\textfont0 = \eightrm\textfont2 = \eightsy
  \left#1\vbox to5.0pt{}\right.\n@space$}}}
\def\tenmath{\tenpoint\fam-1 } % for 10-point in 9-point territory

%%%%%%%%%%%%%%%%%%%%%%%%%%%%%%%%%%%%%%%%%%%%%%%%%%%%%%%%%%%%%%%%%%%%%%%%%%%%
%                    Selecting pages to be printed                         %
%%%%%%%%%%%%%%%%%%%%%%%%%%%%%%%%%%%%%%%%%%%%%%%%%%%%%%%%%%%%%%%%%%%%%%%%%%%%

\let\Shipout=\shipout
\newread\pages \newcount\nextpage \openin\pages=pages
\def\getnextpage{\ifeof\pages\else
 {\endlinechar=-1\read\pages to\next \ifx\next\empty % in this case we
should have eof now \else\global\nextpage=\next\fi}\fi}
\ifeof\pages\else\message{OK, I'll ship only the requested pages!}
 \getnextpage\fi
\def\shipout{\ifeof\pages\let\next=\Shipout
 \else\ifnum\pageno=\nextpage\getnextpage\let\next=\Shipout
  \else\let\next=\Tosspage\fi\fi \next}
\newbox\garbage \def\Tosspage{\deadcycles=0\setbox\garbage=}

%%%%%%%%%%%%%%%%%%%%%%%%%%%%%%%%%%%%%%%%%%%%%%%%%%%%%%%%%%%%%%%%%%%%%%%%%%%%
%                            Two Columns Format                            %
%                         From TeXBook, 1986, D. Knuth                     %
%                                pp. 416-417                               %
%%%%%%%%%%%%%%%%%%%%%%%%%%%%%%%%%%%%%%%%%%%%%%%%%%%%%%%%%%%%%%%%%%%%%%%%%%%%

\ifx\mbrdeuxcol\relax\endinput\else\let\mbrdeuxcol=\relax\fi
%
%    definition of variables
%
% \hoffset = -10 true mm
% \voffset = -5 true mm
% \hsize = 185 true mm
% \vsize = 235 true mm
\newdimen\colwidth
\newdimen\bigcolheight
% \newdimen\separatorwidth
% \separatorwidth = 15 true mm
\newdimen\pagewidth
\pagewidth=\hsize
\newdimen\pageheight
\pageheight=\vsize

\output{\onepageout{\boxmaxdepth=\maxdepth \pagecontents}}
\newbox\partialpage
\newdimen\savesize
%
%    macro \onepageout
%
\def\onepageout#1{\shipout\vbox{\vbox to 0pt{\vskip-22.5pt
                                    \hbox to\pagewidth{\vbox to8.5pt{}
                                          \the\headline}
                                    \vss}\nointerlineskip
                                \vbox to \pageheight{#1}
                                \baselineskip=24pt
                                \hbox to\pagewidth{\the\footline}}
                  \advancepageno
                  \ifnum\outputpenalty>-20000 \else\dosupereject\fi
                 }
%
%     Macro \begindoublecolumns
%
\outer\def\begindoublecolumns{\begingroup
  \pagewidth=\hsize \pageheight=\vsize
  \bigcolheight=\pageheight
  \advance\bigcolheight by\pageheight
  \advance\bigcolheight by\baselineskip
  \colwidth=\pagewidth
  \advance\colwidth by -\separatorwidth
  \divide\colwidth by 2
  \output={\global\setbox\partialpage=\vbox{\unvbox255}}\eject
  \output={\doublecolumnout} \hsize=\colwidth \vsize=\bigcolheight
  \advance\vsize by -2\ht\partialpage}
%
%     Macro \enddoublecolumns
%
\outer\def\enddoublecolumns{\output={\balancecolumns}\eject
  \global\output={\onepageout{\boxmaxdepth=\maxdepth \pagecontents}}
  \global\vsize=\pageheight
  \dimen0=\ht255 \divide\dimen0 by 2
  \ht255=\dimen0
  \endgroup \pagegoal=\vsize}
%
%    Macros called by \begindoublecolumns et \enddoublecolumns
%
\def\doublecolumnout{
  \dimen0=\dp255
  \advance\dimen0 by \ht255
  \divide\dimen0 by 2 \splittopskip=\topskip
  \setbox0=\vsplit255 to\dimen0 \setbox2=\vsplit255 to\dimen0
  \onepageout\pagesofarp
  \global\vsize=\bigcolheight
  \unvbox255 \penalty\outputpenalty\relax}
\def\pagesofarp{\ifvoid\topins\else\unvbox\topins\fi
  \pagesofar
  \ifvoid\footins\else
    \vfill
    \vskip\skip\footins \footnoterule \unvbox\footins \fi}
\def\pagesofar{\unvbox\partialpage
  \wd0=\hsize \wd2=\hsize \hbox to\pagewidth{\box0\hfil\box2}}
\def\balancecolumns{\setbox0=\vbox{\unvbox255} \dimen0=\ht0
  \advance\dimen0 by\topskip \advance\dimen0 by-\baselineskip
  \divide\dimen0 by2 \splittopskip=\topskip
  {\vbadness=10000 \loop \global\setbox3=\copy0
    \global\setbox1=\vsplit3 to\dimen0
    \ifdim\ht3>\dimen0 \global\advance\dimen0 by1pt \repeat}
  \setbox0=\vbox to\dimen0{\unvbox1}
  \setbox2=\vbox to\dimen0{\unvbox3}
  \global\output={\balancingerror}
  \pagesofar}
\newhelp\balerrhelp{Please change the page into one that works}
\def\balancingerror{\errhelp=\balerrhelp
   \errmessage{Page can't be balanced}
   \onepageout{\unvbox255}}
%
%   Modified Plain TeX
%
\catcode`@=11
\def\vfootnote#1{\insert\footins\bgroup
  \interlinepenalty=\interfootnotelinepenalty
  \splittopskip=\ht\strutbox
  \splitmaxdepth=\dp\strutbox \floatingpenalty=20000
  \leftskip=0pt \rightskip=0pt \spaceskip=0pt \xspaceskip=0pt
  \hsize=\pagewidth
  \textindent{#1}\footstrut\futurelet\next\fo@t}
\catcode`@=12

%%%%%%%%%%%%%%%%%%%%%%%%%%%%%%%%%%%%%%%%%%%%%%%%%%%%%%%%%%%%%%%%%%%%%%%%%%%%
%             Bourbaki notations for standard sets: N, Z, Q, ...           %
%%%%%%%%%%%%%%%%%%%%%%%%%%%%%%%%%%%%%%%%%%%%%%%%%%%%%%%%%%%%%%%%%%%%%%%%%%%%

\def\vbar{\mathchoice{\vrule height6.3ptdepth-.5ptwidth.8pt\kern-.8pt}
   {\vrule height6.3ptdepth-.5ptwidth.8pt\kern-.8pt}
   {\vrule height4.1ptdepth-.35ptwidth.6pt\kern-.6pt}
   {\vrule height3.1ptdepth-.25ptwidth.5pt\kern-.5pt}}
\def\fudge{\mathchoice{}{}{\mkern.5mu}{\mkern.8mu}}
\def\bbc#1#2{{\rm \mkern#2mu\vbar\mkern-#2mu#1}}
\def\bbb#1{{\rm I\mkern-3.5mu #1}}
\def\bba#1#2{{\rm #1\mkern-#2mu\fudge #1}}
\def\bb#1{{\count4=`#1 \advance\count4by-64 \ifcase\count4\or\bba A{11.5}\or
   \bbb B\or\bbc C{5}\or\bbb D\or\bbb E\or\bbb F \or\bbc G{5}\or\bbb H\or
   \bbb I\or\bbc J{3}\or\bbb K\or\bbb L \or\bbb M\or\bbb N\or\bbc O{5} \or
   \bbb P\or\bbc Q{5}\or\bbb R\or\bbc S{4.2}\or\bba T{10.5}\or\bbc U{5}\or
%   \bbb P\or\bbc Q{5}\or\bbb R\or\bba S{8}\or\bba T{10.5}\or\bbc U{5}\or
   \bba V{12}\or\bba W{16.5}\or\bba X{11}\or\bba Y{11.7}\or\bba Z{7.5}\fi}}

\def \setN {\bb N} % Natural Integers
\def \setZ {\bb Z} % Rational Integers
\def \setQ {\bb Q} % Rational Numbers
\def \setR {\bb R} % Real Numbers
\def \setC {\bb C} % Complexe Numbers

%%%%%%%%%%%%%%%%%%%%%%%%%%%%%%%%%%%%%%%%%%%%%%%%%%%%%%%%%%%%%%%%%%%%%%%%%%%%
%             BTXMAC: An input file for use of BibTeX with TeX             %
%%%%%%%%%%%%%%%%%%%%%%%%%%%%%%%%%%%%%%%%%%%%%%%%%%%%%%%%%%%%%%%%%%%%%%%%%%%%

\edef\cite{\the\catcode`@}%
\catcode`@ = 11
\let\@oldatcatcode = \cite
\chardef\@letter = 11
\chardef\@other = 12
\def\@innerdef#1#2{\edef#1{\expandafter\noexpand\csname #2\endcsname}}%
\@innerdef\@innernewcount{newcount}%
\@innerdef\@innernewdimen{newdimen}%
\@innerdef\@innernewif{newif}%
\@innerdef\@innernewwrite{newwrite}%
\def\@gobble#1{}%
\ifx\inputlineno\@undefined
   \let\@linenumber = \empty % Pre-3.0.
\else
   \def\@linenumber{\the\inputlineno:\space}%
\fi
\def\@futurenonspacelet#1{\def\cs{#1}%
   \afterassignment\@stepone\let\@nexttoken=
}%
\begingroup % The grouping here avoids stepping on an outside use of `\\'.
\def\\{\global\let\@stoken= }%
\\
\endgroup
\def\@stepone{\expandafter\futurelet\cs\@steptwo}%
\def\@steptwo{\expandafter\ifx\cs\@stoken\let\@@next=\@stepthree
   \else\let\@@next=\@nexttoken\fi \@@next}%
\def\@stepthree{\afterassignment\@stepone\let\@@next= }%
\def\@getoptionalarg#1{%
   \let\@optionaltemp = #1%
   \let\@optionalnext = \relax
   \@futurenonspacelet\@optionalnext\@bracketcheck
}%
\def\@bracketcheck{%
   \ifx [\@optionalnext
      \expandafter\@@getoptionalarg
   \else
      \let\@optionalarg = \empty
      \expandafter\@optionaltemp
   \fi
}%
%
\def\@@getoptionalarg[#1]{%
   \def\@optionalarg{#1}%
   \@optionaltemp
}%
\def\@nnil{\@nil}%
\def\@fornoop#1\@@#2#3{}%
%
\def\@for#1:=#2\do#3{%
   \edef\@fortmp{#2}%
   \ifx\@fortmp\empty \else
      \expandafter\@forloop#2,\@nil,\@nil\@@#1{#3}%
   \fi
}%
%
\def\@forloop#1,#2,#3\@@#4#5{\def#4{#1}\ifx #4\@nnil \else
       #5\def#4{#2}\ifx #4\@nnil \else#5\@iforloop #3\@@#4{#5}\fi\fi
}%
%
\def\@iforloop#1,#2\@@#3#4{\def#3{#1}\ifx #3\@nnil
       \let\@nextwhile=\@fornoop \else
      #4\relax\let\@nextwhile=\@iforloop\fi\@nextwhile#2\@@#3{#4}%
}%
\@innernewif\if@fileexists
%
\def\@testfileexistence{\@getoptionalarg\@finishtestfileexistence}%
\def\@finishtestfileexistence#1{%
   \begingroup
      \def\extension{#1}%
      \immediate\openin0 =
         \ifx\@optionalarg\empty\jobname\else\@optionalarg\fi
         \ifx\extension\empty \else .#1\fi
         \space
      \ifeof 0
         \global\@fileexistsfalse
      \else
         \global\@fileexiststrue
      \fi
      \immediate\closein0
   \endgroup
}%
\def\bibliographystyle#1{%
   \@readauxfile
   \@writeaux{\string\bibstyle{#1}}%
}%
\let\bibstyle = \@gobble
\let\bblfilebasename = \jobname
\def\bibliography#1{%
   \@readauxfile
   \@writeaux{\string\bibdata{#1}}%
   \@testfileexistence[\bblfilebasename]{bbl}%
   \if@fileexists
      % We just output a non-discardable item (the `whatsit' with the
      % \bibdata command).  This means that the glue that will be
      % inserted next (\parskip or \baselineskip, most likely) will be a
      % legal breakpoint.  Most likely, this is after some kind of
      % heading, where we don't want to allow a page break.  So:
      \nobreak
      \@readbblfile
   \fi
}%
\let\bibdata = \@gobble
\def\nocite#1{%
   \@readauxfile
   \@writeaux{\string\citation{#1}}%
}%
%
\@innernewif\if@notfirstcitation
\def\cite{\@getoptionalarg\@cite}%
\def\@cite#1{%
   \let\@citenotetext = \@optionalarg
   \printcitestart
   \nocite{#1}%
   \@notfirstcitationfalse
   \@for \@citation :=#1\do
   {%
      \expandafter\@onecitation\@citation\@@
   }%
   \ifx\empty\@citenotetext\else
      \printcitenote{\@citenotetext}%
   \fi
   \printcitefinish
}%
%
\def\@onecitation#1\@@{%
   \if@notfirstcitation
      \printbetweencitations
   \fi
   %
   \expandafter \ifx \csname\@citelabel{#1}\endcsname \relax
      \if@citewarning
         \message{\@linenumber Undefined citation `#1'.}%
      \fi
      % Give it a dummy definition:
      \expandafter\gdef\csname\@citelabel{#1}\endcsname{%
         {\tt
            \escapechar = -1
            \nobreak\hskip0pt
            \expandafter\string\csname#1\endcsname
            \nobreak\hskip0pt
         }%
      }%
   \fi
   % Now produce the text, whether it was undefined or not.
   \csname\@citelabel{#1}\endcsname
   \@notfirstcitationtrue
}%
\def\@citelabel#1{b@#1}%
\def\@citedef#1#2{\expandafter\gdef\csname\@citelabel{#1}\endcsname{#2}}%
\def\@readbblfile{%
   \ifx\@itemnum\@undefined
      \@innernewcount\@itemnum
   \fi
   %
   \begingroup
      \def\begin##1##2{%
         \setbox0 = \hbox{\biblabelcontents{##2}}%
         \biblabelwidth = \wd0
      }%
      \def\end##1{}% ##1 is `thebibliography' again.
      \@itemnum = 0
      \def\bibitem{\@getoptionalarg\@bibitem}%
      \def\@bibitem{%
         \ifx\@optionalarg\empty
            \expandafter\@numberedbibitem
         \else
            \expandafter\@alphabibitem
         \fi
      }%
      \def\@alphabibitem##1{%
         \expandafter \xdef\csname\@citelabel{##1}\endcsname {\@optionalarg}%
         \ifx\biblabelprecontents\@undefined
            \let\biblabelprecontents = \relax
         \fi
         \ifx\biblabelpostcontents\@undefined
            \let\biblabelpostcontents = \hss
         \fi
         \@finishbibitem{##1}%
      }%
      %
      \def\@numberedbibitem##1{%
         \advance\@itemnum by 1
         \expandafter \xdef\csname\@citelabel{##1}\endcsname{\number\@itemnum}%
         \ifx\biblabelprecontents\@undefined
            \let\biblabelprecontents = \hss
         \fi
         \ifx\biblabelpostcontents\@undefined
            \let\biblabelpostcontents = \relax
         \fi
         \@finishbibitem{##1}%
      }%
      %
      \def\@finishbibitem##1{%
         \biblabelprint{\csname\@citelabel{##1}\endcsname}%
         \@writeaux{\string\@citedef{##1}{\csname\@citelabel{##1}\endcsname}}%
         \ignorespaces
      }%
      \let\em = \bblem
      \let\newblock = \bblnewblock
      \let\sc = \bblsc
      \frenchspacing
      \clubpenalty = 4000 \widowpenalty = 4000
      \tolerance = 10000 \hfuzz = .5pt
      \everypar = {\hangindent = \biblabelwidth
                      \advance\hangindent by \biblabelextraspace}%
      \bblrm
      \parskip = 1.5ex plus .5ex minus .5ex
      \biblabelextraspace = .5em
      \bblhook
      \input \bblfilebasename.bbl
   \endgroup
}%
\@innernewdimen\biblabelwidth
\@innernewdimen\biblabelextraspace
\def\biblabelprint#1{%
   \noindent
   \hbox to \biblabelwidth{%
      \biblabelprecontents
      \biblabelcontents{#1}%
      \biblabelpostcontents
   }%
   \kern\biblabelextraspace
}%
\def\biblabelcontents#1{{\bblrm [#1]}}%
\def\bblrm{\rm}%
\def\bblem{\it}%
\def\bblsc{\ifx\@scfont\@undefined
              \font\@scfont = cmcsc10
           \fi
           \@scfont
}%
\def\bblnewblock{\hskip .11em plus .33em minus .07em }%
\let\bblhook = \empty
\def\printcitestart{[}%         left bracket
\def\printcitefinish{]}%        right bracket
\def\printbetweencitations{, }% comma, space
\def\printcitenote#1{, #1}%     comma, space, note (if it exists)
\let\citation = \@gobble
\@innernewcount\@numparams
\def\newcommand#1{%
   \def\@commandname{#1}%
   \@getoptionalarg\@continuenewcommand
}%
\def\@continuenewcommand{%
   \@numparams = \ifx\@optionalarg\empty 0\else\@optionalarg \fi \relax
   \@newcommand
}%
\def\@newcommand#1{%
   \def\@startdef{\expandafter\edef\@commandname}%
   \ifnum\@numparams=0
      \let\@paramdef = \empty
   \else
      \ifnum\@numparams>9
         \errmessage{\the\@numparams\space is too many parameters}%
      \else
         \ifnum\@numparams<0
            \errmessage{\the\@numparams\space is too few parameters}%
         \else
            \edef\@paramdef{%
               % This is disgusting, but \loop doesn't work inside \edef,
               % because \body isn't defined.
               \ifcase\@numparams
                  \empty  No arguments.
               \or ####1%
               \or ####1####2%
               \or ####1####2####3%
               \or ####1####2####3####4%
               \or ####1####2####3####4####5%
               \or ####1####2####3####4####5####6%
               \or ####1####2####3####4####5####6####7%
               \or ####1####2####3####4####5####6####7####8%
               \or ####1####2####3####4####5####6####7####8####9%
               \fi
            }%
         \fi
      \fi
   \fi
   \expandafter\@startdef\@paramdef{#1}%
}%
\def\@readauxfile{%
   \if@auxfiledone \else % remember: \@auxfiledonetrue if \noauxfile is defined
      \global\@auxfiledonetrue
      \@testfileexistence{aux}%
      \if@fileexists
         \begingroup
            \endlinechar = -1
            \catcode`@ = 11
            \input \jobname.aux
         \endgroup
      \else
         \message{\@undefinedmessage}%
         \global\@citewarningfalse
      \fi
      \immediate\openout\@auxfile = \jobname.aux
   \fi
}%
\newif\if@auxfiledone
\ifx\noauxfile\@undefined \else \@auxfiledonetrue\fi
\@innernewwrite\@auxfile
\def\@writeaux#1{\ifx\noauxfile\@undefined \write\@auxfile{#1}\fi}%
\ifx\@undefinedmessage\@undefined
   \def\@undefinedmessage{No .aux file; I won't give you warnings about
                          undefined citations.}%
\fi
\@innernewif\if@citewarning
\ifx\noauxfile\@undefined \@citewarningtrue\fi
\catcode`@ = \@oldatcatcode

%%%%%%%%%%%%%%%%%%%%%%%%%%%%%%%%%%%%%%%%%%%%%%%%%%%%%%%%%%%%%%%%%%%%%%%%%%%%
%                   Macro to allow block comments in TeX                   %
%                       J.C. Alexander,  May, 1986                         %
%%%%%%%%%%%%%%%%%%%%%%%%%%%%%%%%%%%%%%%%%%%%%%%%%%%%%%%%%%%%%%%%%%%%%%%%%%%%

\edef\Saveatcatcode{\the\catcode`\@}
\catcode`\@=11
\def\newcodes@{\edef\S@veslashcatcode{\the\catcode`\\}
               \edef\S@velbraccatcode{\the\catcode`\{}
               \edef\S@verbraccatcode{\the\catcode`\}}
               \edef\S@venumsgcatcode{\the\catcode`\#}
               \edef\S@veperctcatcode{\the\catcode`\%}
\catcode`\\=12 \catcode`\{=12 \catcode`\}=12 \catcode`\#=12
\catcode`\%=12\relax}
\def\oldcodes@{\catcode`\\=\S@veslashcatcode
               \catcode`\{=\S@velbraccatcode
               \catcode`\}=\S@verbraccatcode
               \catcode`\#=\S@venumsgcatcode
               \catcode`\%=\S@veperctcatcode
               \relax}
\def\begincomment{\newcodes@\endlinechar=10 \comment@}
{\lccode`\!=`\\
\lowercase{\gdef\comment@#1^^J{\comment@@#1!endcomment\comment@@@}%
\gdef\comment@@#1!endcomment{\futurelet\next\comment@@@}%
\gdef\comment@@@#1\comment@@@{\ifx\next\comment@@@\let
\next=\comment@\else\def\next{\oldcodes@\endlinechar=`\^^M\relax}%
\fi\next}}}
\catcode`\@=\Saveatcatcode

%%%%%%%%%%%%%%%%%%%%%%%%%%%%%%%%%%%%%%%%%%%%%%%%%%%%%%%%%%%%%%%%%%%%%%%%%%%%
%                                   Misc.                                  %
%%%%%%%%%%%%%%%%%%%%%%%%%%%%%%%%%%%%%%%%%%%%%%%%%%%%%%%%%%%%%%%%%%%%%%%%%%%%

\def \BA {{\bf A}}
\def \BB {{\bf B}}
\def \BC {{\bf C}}
\def \BE {{\bf E}}
\def \BF {{\bf F}}
\def \BG {{\bf G}}
\def \BH {{\bf H}}
\def \BI {{\bf I}}
\def \BJ {{\bf J}}
\def \Bk {{\bf k}}
\def \BK {{\bf K}}
\def \BL {{\bf L}}
\def \BM {{\bf M}}
\def \BN {{\bf N}}
\def \BQ {{\bf Q}}
\def \BP {{\bf P}}
\def \BQ {{\bf Q}}
\def \BR {{\bf R}}
\def \BS {{\bf S}}
\def \BV {{\bf V}}

\def \Ga   {{\eufm a}}
\def \CA   {{\eusm A}}
\def \Gb   {{\eufm b}}
\def \CB   {{\eusm B}}
\def \CC   {{\eusm C}}
\def \CD   {{\eusm D}}
\def \CE   {{\eusm E}}
\def \CF   {{\eusm F}}
\def \CG   {{\eusm G}}
\def \CH   {{\eusm H}}
\def \CI   {{\eusm I}}
\def \GI   {{\eufm I}}
\def \CJ   {{\eusm J}}
\def \CK   {{\eusm K}}
\def \CL   {{\eusm L}}
\def \CM   {{\eusm M}}
\def \Cm   {{\eurm m}}
\def \Gm   {{\eufm m}}
\def \CN   {{\eusm N}}
\def \Gp   {{\eufm p}}
\def \CP   {{\eusm P}}
\def \Gq   {{\eufm q}}
\def \CQ   {{\eusm Q}}
\def \CR   {{\eusm R}}
\def \CS   {{\eusm S}}
\def \CT   {{\eusm T}}
\def \CU   {{\eusm U}}
\def \CV   {{\eusm V}}
\def \CW   {{\eusm W}}
\def \CX   {{\eusm X}}
\def \CXeb {{\eusm X}_{\rm eb}}
\def \CY   {{\eusm Y}}
\def \CZ   {{\eusm Z}}

% Footnote symbols
\def \fndag       {\hbox{${}^{{}^{\dag}}$}}
\def \fnast       {\hbox{${}^{{}^\ast}$}}
\def \fndast      {\hbox{${}^{{}^{\ast\ast}}$}}
\def \fntast      {\hbox{${}^{{}^{\ast\ast\ast}}$}}
\def \fnone       {\hbox{${}^{{}^1}$}}
\def \fntwo       {\hbox{${}^{{}^2}$}}
\def \fnthree     {\hbox{${}^{{}^3}$}}
\def \fnfour      {\hbox{${}^{{}^4}$}}
\def \fnfive      {\hbox{${}^{{}^5}$}}
\def \fnsix       {\hbox{${}^{{}^6}$}}
\def \fnseven     {\hbox{${}^{{}^7}$}}
\def \fneight     {\hbox{${}^{{}^8}$}}
\def \fnnine      {\hbox{${}^{{}^9}$}}
\def \fnten       {\hbox{${}^{{}^10}$}}
%%%%%%%%%%%%%%%%%%%%%%%%%%%%%%%%%%%%%%%%%%%%%%%%%%%%%%%%%%%%%%%%%%%%%%%%%%%%
%%%%%%%%%%%%%%%%%%%%%%%%%%%%%%%%%%%%%%%%%%%%%%%%%%%%%%%%%%%%%%%%%%%%%%%%%%%%




